\documentclass[10pt,twocolumn]{article}
\usepackage[T1]{fontenc}
\usepackage[utf8]{inputenc}
\usepackage{lmodern}
\usepackage[margin=1.8cm,top=2.4cm,bottom=2cm,columnsep=0.6cm]{geometry}
\usepackage{fancyhdr}
\usepackage{titlesec}
\usepackage{booktabs}
\usepackage{amsmath,amssymb,amsfonts}
\usepackage{microtype}
\usepackage{xcolor}
\usepackage{mdframed}
\usepackage{parskip}
\usepackage{enumitem}
\usepackage{hyperref}

\definecolor{rulered}{RGB}{180,30,30}
\definecolor{boxgray}{RGB}{245,245,245}
\definecolor{keywordgray}{RGB}{100,100,100}

\pagestyle{fancy}
\fancyhf{}
\fancyhead[L]{\textsc{\small Claude Mag}}
\fancyhead[C]{\small\textit{Machine Learning Research Digest}}
\fancyhead[R]{\small 2026-10-07}
\fancyfoot[C]{\small\thepage}
\renewcommand{\headrulewidth}{0.8pt}

\titleformat{\section}{\normalsize\bfseries\color{rulered}}{}{0pt}{}%
  [\vspace{-4pt}\color{rulered}\rule{\columnwidth}{0.4pt}]
\titlespacing{\section}{0pt}{8pt}{4pt}

\hypersetup{colorlinks=true,urlcolor=rulered,linkcolor=black}

\begin{document}
\setlength{\parindent}{0pt}
\setlength{\parskip}{4pt}

{\small\color{keywordgray}\textbf{Keywords:}
  web agents \textbullet{}
  conformal prediction \textbullet{}
  reinforcement learning \textbullet{}
  test-time scaling}\par\vspace{4pt}

{\LARGE\bfseries\raggedright
  CLIFT: Conformal Self-Verification for Web Agent Training and
  Test-Time Scaling}\par\vspace{4pt}

{\large\itshape\raggedright
  A Conformally Certified URL-Aware Verifier Replaces Costly Frontier-LM
  Supervision with Dense Per-Step Signals and Enables Reliable Test-Time
  Scaling Without External Oracles}\par\vspace{6pt}

{\small\textbf{By} Yifan Zhang, Yutong Dai, Viraj Prabhu, Zhiyuan Hu,
  Ran Xu \& Zeyuan Chen
  (Georgia Institute of Technology; University of Texas at Austin)}\par\vspace{2pt}

{\small\color{keywordgray}arXiv:2610.06829 \textbullet{} October 2026}
\par\vspace{2pt}\noindent{\color{rulered}\rule{\columnwidth}{0.6pt}}\vspace{6pt}

Web agent training faces a fundamental credit-assignment problem: binary
task-success rewards are sparse, while frontier-LM judges providing denser
step-level feedback are expensive and unavailable at deployment. CLIFT solves
both at once with a conformally calibrated URL-aware verifier that supplies
non-negative per-step evidence during training and drives reliable test-time
scaling without querying any external oracle.

\begin{mdframed}[backgroundcolor=boxgray,linecolor=rulered,linewidth=1pt,
    innerleftmargin=6pt,innerrightmargin=6pt,
    innertopmargin=6pt,innerbottommargin=6pt]
\textbf{\small AT A GLANCE}\par\vspace{4pt}
\begin{description}[leftmargin=0pt,labelsep=4pt,itemsep=2pt,parsep=0pt,
    font=\small\bfseries]
  \item[Problem:] Sparse binary task rewards leave most intermediate actions
    unscored; frontier-LM judges are too expensive for every training step.
  \item[Why it matters:] Dense, calibrated per-step signals accelerate web
    agent training and enable reliable test-time compute scaling.
  \item[Method:] Natural-language verification questions grounded in browser
    state, calibrated by conformal prediction; same verifier reused at
    inference for Conformal Test-time Scaling.
  \item[Result:] Higher task success on WebArena and VisualWebArena; 8--12 pp
    gain from test-time scaling with $N{=}8$ trajectories, no oracle needed.
  \item[Evidence:] WebArena and VisualWebArena benchmarks; ablations on
    calibration, question quality, and verifier model size.
\end{description}
\end{mdframed}

\section{The Big Picture}

Reinforcement learning from task-completion rewards is the dominant paradigm
for training web agents, but it suffers a fundamental supervision bottleneck:
a twenty-step trajectory that ultimately succeeds returns the same sparse
positive signal regardless of which intermediate actions were actually useful.
Frontier language model judges can provide step-level feedback, but querying
GPT-4o at every step during training is cost-prohibitive and creates
deployment-time dependencies on proprietary APIs.

CLIFT (Conformal Learning for Instruction Following at Test time) addresses
both problems simultaneously. During training, it replaces frontier-LM
supervision with a small verifier whose outputs are calibrated by conformal
prediction, guaranteeing non-negative evidence in expectation. At inference,
the same frozen verifier drives Conformal Test-time Scaling (CTS), selecting
among candidate trajectories without any external oracle.

\section{Why Existing Approaches Fall Short}

\textbf{Binary rewards are sparse.} An agent that completes a task in twenty
steps receives no information about which individual steps contributed. This
makes gradient estimation difficult and exploration slow.

\textbf{LM judges are expensive and API-dependent.} Querying a frontier model
at every step dramatically increases training cost and creates an availability
dependency that is absent at deployment.

\textbf{Heuristic verifiers lack calibration.} Rule-based verifiers checking
URL patterns or DOM state can fire incorrectly; without calibration their
signals may actively mislead training.

\textbf{Test-time scaling is uncertified.} Best-of-$N$ and chain-of-thought
approaches for agents offer no statistical guarantee that more compute
improves reliability.

\section{Core Mechanism}

CLIFT introduces two complementary components:

\textbf{Conformal Self-Verification (training):} Natural-language verification
questions (e.g., ``Is the agent on the correct URL?'', ``Has the required
form field been filled?'') are generated from task descriptions and grounded
in the current browser state (URL and DOM excerpt). A small verifier answers
these questions and assigns per-step scores. Conformal prediction calibrates
the verifier's confidence so it fires only when it can certify its answer at
a pre-specified error rate---yielding a conformally certified, URL-aware
verifier with a formal coverage guarantee.

\textbf{Conformal Test-time Scaling (CTS, inference):} The same frozen
verifier bank scores all candidate trajectories sampled by the agent. CTS
selects the trajectory with the highest cumulative conformal evidence,
enabling reliable test-time scaling with no frontier-model dependency.

\section{Technical Formulation}

Let $s_t = (u_t, d_t)$ be the browser state at step $t$ (URL and DOM excerpt),
$a_t$ the agent's action, and $Q = \{q_1, \ldots, q_K\}$ a bank of
verification questions derived from the task description.

The verifier model $v_\phi$ scores each question-state pair:
\[
p_{k,t} = v_\phi(q_k \mid s_t, a_t) \in [0, 1]
\]

Conformal calibration with coverage $1 - \alpha$ produces a threshold $\tau$
such that:
\[
\Pr[\text{verifier incorrect}] \leq \alpha
\]

The per-step evidence signal used during training is:
\[
r_t = \sum_{k=1}^{K} \mathbf{1}[p_{k,t} \geq \tau] \cdot w_k
\]
where $w_k$ are task-specific weights. This signal is always non-negative
(evidence of progress, never a penalty), so it supplements rather than
replaces the sparse task-completion reward.

For CTS at inference, given $N$ sampled trajectories
$\{\tau^{(i)}\}_{i=1}^N$, the selected trajectory is:
\[
\hat{\tau} = \arg\max_i \sum_t r_t^{(i)}
\]

\section{Learning or Inference Procedure}

\textbf{Training:}
\begin{enumerate}[itemsep=2pt,parsep=0pt]
  \item Generate a bank of verification questions per task type.
  \item Calibrate the verifier on a held-out set to determine $\tau$
    at desired coverage $1 - \alpha$.
  \item Run RL with task-completion reward augmented by conformal
    per-step evidence $r_t$.
\end{enumerate}

\textbf{Test-time scaling:}
\begin{enumerate}[itemsep=2pt,parsep=0pt]
  \item Sample $N$ trajectories from the policy.
  \item Score each step with the frozen conformal verifier.
  \item Return the trajectory with the highest cumulative evidence.
\end{enumerate}

\section{What the Guarantee Says}

By construction of conformal prediction, the verifier provides a marginal
coverage guarantee: the probability that the verifier incorrectly certifies an
incorrect action as correct is at most $\alpha$ (typically $\alpha = 0.1$).
This ensures per-step signals are non-negative in expectation, preserving
policy improvement under the RL update. The CTS selection guarantee is weaker,
but conformal calibration ensures the selected trajectory is at least as
likely to be correct as random selection would predict.

\section{Experimental Findings}

\begin{itemize}[itemsep=2pt,parsep=0pt]
  \item \textbf{WebArena and VisualWebArena:} CLIFT-augmented RL achieves
    significantly higher task success than RL with binary reward alone.
  \item \textbf{Evidence density:} The conformal verifier fires at roughly
    40--60\% of steps, providing substantially denser credit assignment than
    binary rewards.
  \item \textbf{CTS efficiency:} $N{=}8$ trajectories with CTS selection
    improves success rate by \textbf{8--12 percentage points} over
    single-trajectory evaluation, with no frontier-model query.
  \item \textbf{Calibration quality:} The conformally calibrated verifier
    achieves coverage close to the nominal $1 - \alpha$ across all task types.
\end{itemize}

\section{Ablations and Interpretation}

\textbf{Calibration vs.\ uncalibrated verifier:} Using the raw verifier
without conformal calibration produces noisy per-step signals that can harm
training; calibration is critical to the method's benefit.

\textbf{Question quality:} Generic verification questions degrade performance;
the quality of the task-specific question bank governs verifier coverage.

\textbf{Number of CTS samples:} Performance improves with more trajectories
up to $N{=}16$; beyond that, gains are marginal, suggesting the verifier
reaches its discrimination ceiling.

\textbf{Verifier model size:} A 1B-parameter verifier delivers comparable
quality to a 7B verifier for structured URL-aware verification, indicating
that small specialized models are sufficient.

\section*{Reference}

Yifan Zhang, Yutong Dai, Viraj Prabhu, Zhiyuan Hu, Ran Xu \& Zeyuan Chen.
\textbf{CLIFT: Conformal Self-Verification for Web Agent Training and
Test-Time Scaling}.
arXiv:2610.06829, October 2026.\\
\url{https://arxiv.org/abs/2610.06829}

\end{document}
